% !TEX root = ../main.tex
\chapter{Why this exists, and what it has found}\label{ch:why}
\section{The question}\label{sec:why-question}
Fungi grow as a network of fine threads, the mycelium, through soil and wood. Electrodes
pushed into that network pick up a small, slowly changing voltage. The project behind this
app records that voltage on many electrodes at once, for days at a time, and asks whether
the activity contains \emph{recurring patterns}: shapes that come back again and again, on
one electrode or across several.

\begin{plainterms}[What a recording looks like]
Picture a heart monitor that runs for a month instead of a minute, with sixteen leads
instead of one, and a trace that moves so slowly that a single ``beat'' can last minutes.
Somewhere in that month of wiggles there may be shapes that repeat. Finding them is the
job.
\end{plainterms}

Three things make the job hard.
\begin{itemize}
  \item \textbf{Scale.} One recording in the project's library, \code{M2_aug_concat_fs1.mat},
    holds 16 channels of about 721 hours each, one sample a second. Nobody can look at all
    of that by eye, let alone compare it with the next recording.
  \item \textbf{Judgement.} Whether a wiggle is ``interesting'' is, in the end, a person's
    call. The only labels the project had to start with were made by one person, by hand
    (\code{docs/PIPELINE_PRD.md}, ``Problem Statement'').
  \item \textbf{Chance.} A long, noisy trace will contain look-alike shapes by accident. A
    claim that a pattern recurs means something only if it recurs \emph{more often than it
    would by chance}.
\end{itemize}

Before the app, each analysis technique lived in its own script. The product
specification describes the cost plainly: a promising result was found and then lost,
because the settings and the order of steps that produced it were never saved in a form
that could be run again (\code{docs/PIPELINE_PRD.md}, ``Problem Statement'').
\unverified{The researcher built the app so that the thesis claim, that recurring patterns
in fungal electrical recordings can be detected and characterised automatically, rests on
evidence that can be regenerated rather than on scattered plots.}

\section{What the app does about it}\label{sec:why-does}
The app gives that work one shape, which it draws across the top of every page:
\begin{center}\sffamily
\textbf{Explore}: find a span \enskip\ensuremath{\to}\enskip
\textbf{Analyse}: run a chain over it \enskip\ensuremath{\to}\enskip
\textbf{Review}: judge what it found
\end{center}

\begin{itemize}
  \item In \textbf{Explore} you look at a recording: every channel at once, then one channel
    at any zoom, then one candidate shape. You pick a stretch of time, a \term{span}{span}.
  \item In \textbf{Analyse} you build a \term{chain}{chain} of \term{block}{blocks}: small,
    named analysis steps, each doing one thing (remove the slow drift, find the spikes, group
    similar shapes). You run the chain over your span and see every step's result on the
    same time axis.
  \item In \textbf{Review} you judge, one at a time, what the machine proposed, with one of
    five fixed \term{verdict}{verdicts}.
\end{itemize}
Around that path, \textbf{Discovery} runs a saved chain across many channels and checks it
against chance, \textbf{Library} keeps the shapes that survived your judgement as
\term{family}{families} you can search for again, and \textbf{Models} trains classifiers on
what Analyse produced. \Cref{ch:core-ideas} explains each of these words.

\begin{plainterms}[A recipe book, not a pile of notes]
A chain is like a written recipe. Anyone can follow it again, on the same ingredients or on
new ones, and get the same dish. The old scripts were like a cook's memory: the meal was
good, but nobody could say exactly how it was made.
\end{plainterms}

\section{What it has found so far}\label{sec:why-found}
Short answer: the instrument is built, and the research questions are only now being
answered with it. The project's specification names six research questions
(\code{docs/PIPELINE_PRD.md}, ``How the research questions map onto the build''); each has
its own living file in \code{docs/rq_roundA/} saying how answerable it is. As those files
stand:
\begin{itemize}
  \item \textbf{Question 1}, whether labels made by clustering train a better classifier than
    labels made by hand: answerable for its first two tests; a first paired run exists, and
    the main result has been re-scoped (\code{docs/rq_roundA/RQ1-cluster-vs-manual-labels.md}).
  \item \textbf{Question 2}, whether a shape you have checked, used as a search seed, finds
    more shapes you also accept: answerable in the app, not yet answered
    (\code{RQ2-exemplar-seeded-search.md}).
  \item \textbf{Question 3}, whether a motif stays ``the same'' when stretched or squashed in
    time: answerable in the app, not yet answered; the first read-outs used placeholder
    verdicts, so they show the instrument works, not what the answer is
    (\code{RQ3-identity-under-scale.md}).
  \item \textbf{Question 4}, whether splitting the signal into frequency bands and turning it
    into letters finds regions the plain signal misses: answerable in the app
    (\code{RQ4-bands-and-symbols.md}).
  \item \textbf{Question 5}, where human and machine judgement disagree: the instrument is
    built, but the human verdicts it needs have not yet been given in Review
    (\code{RQ5-human-vs-machine-divergence.md}).
  \item \textbf{Question 6}, whether a pattern recurs across channels once crossed wires and
    travelling signals are taken out: answerable on families whose members are long enough,
    not yet answered (\code{RQ6-recurrence-across-channels.md}).
\end{itemize}

One finding recurs across these files and is worth knowing before you start:
\emph{at the default settings, the runs measured so far did not beat chance}. The readiness
report of 2 October 2026 states that nothing it ran came out ahead of its surrogate, and the
band runs of Question~4 found fewer regions on the real signal than their chance version
predicted (\code{docs/RESEARCH_READINESS.md}, ``Three things that bear on more than one
question''; \code{docs/rq_roundA/RQ4-bands-and-symbols.md}). That is a statement about
untuned defaults, not about fungi. It is the reason every result in the app is shown next to
its \term{null}{null}.

\section{Why it is built the way it is}\label{sec:why-built}
Four choices shape everything else in the app. Each answers one of the three difficulties
above.

\subsection*{Every run can be repeated exactly}
A chain is saved as a recipe: the blocks, their order and every setting. The recipe is
turned into a short fingerprint (a hash), and every result, cached step, saved plot and
Library link carries the fingerprint of the recipe that made it
(\code{docs/PIPELINE_PRD.md}, ``Solution''). If two results carry the same fingerprint, the
same recipe made them.

\begin{plainterms}[Fingerprints]
A recipe's fingerprint works like a barcode on a jar of jam. Two jars with the same barcode
came from the same recipe; change one ingredient and the barcode changes too.
\end{plainterms}

\subsection*{The machine proposes; a person decides}
What a block finds is a \term{detection}{detection}: a machine's proposal. What a person
marks is an \term{annotation}{annotation}. The two are kept in separate tables in the
database, so that no machine output can ever be mistaken for a human judgement, or the
reverse. A person's verdict on a detection goes into a third table of its own, and nothing
enters the Library without a person's say-so. This is the project's rule 5
(\code{CLAUDE.md}); \cref{sec:core-ideas-detections} explains why it matters.

\begin{plainterms}[Why keep them apart]
Think of an exam marked by a teacher. If the pupils' answers and the teacher's marks were
written in the same column, nobody could later ask ``how often was the pupil right?''
Keeping the machine's answers and the person's marks in separate columns keeps that question
open.
\end{plainterms}

\subsection*{Every result is checked against chance}
The app can build a fake version of the recording that keeps its general character but
scrambles any real pattern (a \term{surrogate}{surrogate}), run the same chain on it, and
show how much the chain finds there. A result counts only if the real recording beats its
fake copies. \Cref{sec:core-ideas-nulls} explains how, and \cref{app:nulls} gives the detail.

\subsection*{The same code runs on a laptop and on a cluster}
The analysis code underneath the app (\code{Working/} and \code{Adapters/}) knows nothing
about the browser: it imports no web or drawing library, and a test enforces this
(\code{tests/test_import_boundaries.py}). So the same recipe that runs behind the Analyse
page can run unattended, many copies at once, on a university computing cluster, and its
results can be brought back into the app. Long searches are meant to go there; the
\term{job}{Jobs} page tracks them.

\begin{underthehood}[Where to read more]
Purpose and design: \code{docs/PIPELINE_PRD.md} Part~1 (``Problem Statement'',
``Solution'', ``Implementation Decisions''); the dates in it are a record of the original
plan, not a live schedule. The six questions: \code{docs/rq_roundA/README.md} and one file
per question. The measured state of each question on 2 October 2026:
\code{docs/RESEARCH_READINESS.md}. How the web UI was wired to the core:
\code{docs/WIRING_PLAN.md}. \Cref{ch:architecture} describes the build.
\end{underthehood}
